The integration of Large Language Models (LLMs) as conversational learning assistants and virtual tutors introduces significant challenges for software engineering. Unlike conventional software systems governed by deterministic logic, generative language models operate via probabilistic token sampling. This non-deterministic nature prevents the use of standard boolean assertion testing and represents a direct manifestation of the oracle problem. In educational environments, such stochastic behavior often results in plausible yet incorrect hallucinations, security vulnerabilities against adversarial prompt injection, and passive solutionism that undermines student critical thinking.

This Bachelor's Thesis proposes and validates a formal, systematic testing methodology designed to evaluate the quality, reliability, and pedagogical effectiveness of LLM-based conversational systems. Grounded in the software product quality attributes of ISO/IEC 25010:2023, along with established pedagogical frameworks such as Vygotsky's Zone of Proximal Development and Hattie and Timperley's formative feedback model, the proposed framework establishes seven orthogonal evaluation dimensions. To eliminate the central tendency bias typically found in odd Likert scales, an even four-level analytical rubric ($0$ to $3$) with objective behavioral descriptors is introduced, supported by aggregated quantitative metrics and a Safety-First Educational Quality Index ($IQE$).

The methodology was experimentally validated using a test suite of 42 structured test cases evaluated across three distinct chatbot profiles: a baseline assistant lacking specific pedagogical instructions, a direct expository tutor, and a socratic scaffolding tutor. Empirical results demonstrate that an unaligned baseline model exhibits a 33.3\,\% hallucination rate on deceptive queries and a 9.5\,\% critical safety failure rate, whereas enforcing pedagogical system constraints and socratic prompting achieves full educational suitability. Furthermore, inter-annotator agreement measured via Cohen's kappa ($\kappa = 0.931$) confirms the high reliability and reproducibility of the proposed evaluation process.

\keywords{Large Language Models, Software Testing, Educational Chatbots, Software Quality, Hallucinations, Socratic Scaffolding, ISO/IEC 25010, Artificial Intelligence in Education}
